\section{Introduction}
\label{sec:intro}

This article describes a single, machine-verified idea and the discipline that keeps it honest.
The idea is a physical prohibition: an engine that performs the nontrivial work of descent
\emph{forever}, losing nothing, cannot exist. In the arithmetic of the natural numbers this is
Fermat's method of infinite descent; we restate it for the centres $m$ of the Euclidean pairs
$(6m-1,\,6m+1)$ and call the forbidden object \emph{Euclid's perpetual engine}. The prohibition
itself --- there is no infinitely descending chain of clean Euclidean steps --- is proved on the
bare Lean~4 kernel, without \code{mathlib} and without any axiom beyond the standard ones
(\cref{sec:engine}).

The thesis of the programme is that a surprising number of hard questions are shadows of this one
prohibition cast on different objects. Where an object cannot host a perpetual engine, a
structural ``no deviation'' theorem holds \emph{unconditionally}; what remains open is always the
same kind of thing --- the tie between the abstract structure and a concrete analytic object that
mathematics has not yet supplied. We carry this reading through the twin-prime conjecture (the
spine of the article), the Riemann hypothesis, $\mathrm{P}$ versus $\mathrm{NP}$, Yang--Mills,
Hodge, Navier--Stokes, Birch--Swinnerton-Dyer, and further through Mersenne primes, an arithmetic
zoo of prime patterns, the geometry of the descent graph, and the Collatz problem.

\paragraph{What is proved, and what is not.}
We are emphatic on this point, because the subject invites overclaiming and we wish to invite
none. The twin-prime conjecture is \emph{not} proved here. What is proved, and machine-checked
without \code{sorry}, is a \emph{reduction}: the entire twin branch is brought down to one
explicit open node --- a finite semantic key that must resolve certain genealogy collisions
(\cref{sec:reduction}). That node cannot be closed from inside the system: closing it would amount
to running the very engine whose impossibility is the foundation. We therefore accept it
\emph{from outside}, as a single axiom we call the \emph{first cause} (\cref{sec:firstcause}),
kept in one quarantined module. Our main theorem, \emph{higher-energy incompatibility}, then makes
the situation precise: the infinitude of the twins can be neither proved nor refuted within the
system without a perpetual engine, so a finite observer can at most \emph{accept} the first cause,
never derive it.

\paragraph{The honesty discipline.}
Three colours run through the development and we name them plainly: \green{}~a statement is
machine-proved under the standard axioms (this includes a \emph{green conditional theorem}, whose
hypothesis is a named input rather than the axiom); \yellow{}~a statement is \emph{axiom-tainted},
i.e.\ it depends on the first cause; \red{}~an object is genuinely open. A separate verifier walks
over every declaration of the repository, collects the axioms each one depends on, and enforces
two invariants that never change: the only non-standard axiom in the entire development is
\axiom{}, and the only \code{sorry}-tainted statement is the twin-prime conjecture itself. At the
time of writing the verifier reports exactly one non-standard axiom, exactly one \code{sorry}, and
$47$ axiom-tainted declarations; a green module cannot be forged, and the axiom-tracker leaves a
decree nowhere to hide behind a green label. Four times during the work an internal audit found an
empty wrapper before it reached the showcase (we record these vacuities where they occurred);
that they were caught by machine is the best reason to trust the parts that passed.

\paragraph{Reproducibility.}
Everything in this article is a reflection of Lean~4 source
code~\cite{lean4,mathlib,euclidspath-repo}. Each formal statement below is tagged with the name of
its machine-verified declaration, printed in \code{typewriter} font. The complete source is public
at \url{https://github.com/elamaunt/Euclids-path}; the identical content is browsable as a
bilingual site at \url{https://elamaunt.github.io/Euclids-path/}. Building with
\code{lake build} checks the whole development, and \srcpath{lake env lean scripts/VerifyAll.lean}
re-establishes the two honesty invariants and fails loudly if either is violated.

\paragraph{Plan.}
\cref{sec:engine} proves the impossibility of the engine and its irreversibility.
\cref{sec:carrier} records the carrier of two and the conservation laws that make the pairs
``held apart''. \cref{sec:reduction} reduces the twin conjecture to a single node. \cref{sec:firstcause}
introduces the first cause and proves the main theorem. \cref{sec:riemann,sec:pnp,sec:ymhodge,sec:ns}
read the Riemann hypothesis, $\mathrm{P}$ versus $\mathrm{NP}$, Yang--Mills and Hodge, and
Navier--Stokes through the engine. \cref{sec:zoo,sec:geometry,sec:bsd} treat the Mersenne branch
and the arithmetic zoo, the geometry of the path, and Birch--Swinnerton-Dyer together with several
open neighbours. \cref{sec:collatz} treats the Collatz problem, including a decree that was taken
and then withdrawn when its law was machine-refuted. \cref{sec:coda} is a synthesis;
\cref{sec:status} gives the honest global status and the verification procedure; \cref{sec:figures}
specifies the generation algorithm of each figure.
